% Responsable: P2  (versión completa en extendido/03_parser.tex)
\section{Implementación del parser}
\label{sec:parser}

El parser (\texttt{spanish/parser.py}) es descendente recursivo: cada no terminal es un método
(\texttt{sentence}, \texttt{noun\_phrase}, \texttt{verb\_phrase}…). Las reglas con
$\varepsilon$ son un \texttt{if} y las reglas $X' \to \alpha X' \mid \varepsilon$ son un
\texttt{while}, ambos guiados por el conjunto FIRST de la palabra actual. Una oración se acepta
solo si se construye el árbol y no sobran palabras.

\paragraph{Palabras ambiguas.} El lexer entrega todas las lecturas de cada palabra y el parser
toma la que pide su posición: en \emph{El gato duerme bajo la mesa}, \emph{bajo} es preposición;
en \emph{El hombre bajo come}, adjetivo. Sin retroceso, esto falla en un caso que la gramática sí
genera: en \emph{El gato bajo la mesa duerme} el parser toma \emph{bajo} como adjetivo y luego no
puede seguir. Resolverlo requiere retroceso o un parser de gráfico.

\paragraph{Pipeline.} \texttt{spanish/pipeline.py} encadena lexer, parser, analizador de
concordancia y visualización; si una etapa falla, la siguiente no se ejecuta. Cada error indica
la palabra y la causa:
\begin{center}
  \small
  \resizebox{\linewidth}{!}{%
  \begin{tabular}{ll}
    \toprule
    Oración & Salida \\
    \midrule
    \emph{El perro corre en el parque.} & sintáctico, palabra 3: `corre' no está en la gramática \\
    \emph{Come la manzana.} & sintáctico, palabra 1: se esperaba un pronombre, determinante o posesivo \\
    \emph{Los perros come.} & concordancia, palabra 3: el verbo `come' (singular) no concuerda con `perros' (plural) \\
    \bottomrule
  \end{tabular}}
\end{center}
El resultado incluye además las etiquetas, el sujeto, el verbo, el objeto y la unión de la
frase preposicional, que se sacan del árbol. Hay 220 pruebas automáticas; una de ellas pasa las
73 oraciones del conjunto de prueba por el pipeline completo.
